\title{QLoRA: Efficient Finetuning of Quantized LLMs}

\begin{abstract}
We introduce \textsc{QLoRA}, a memory-efficient finetuning method that enables the finetuning of models with up to 65B parameters on a single 48GB GPU, all while retaining the full effectiveness of 16-bit task performance. \textsc{QLoRA} achieves this by computing gradients through a fixed, 4-bit quantized pretrained language model and directing them into Low Rank Adapters (LoRA). Our leading model series, named \textbf{Guanaco}, establishes a new standard among open-source models by surpassing previous best performances on the Vicuna benchmark, reaching 99.3\% of ChatGPT’s capability with just 24 hours of single-GPU finetuning. \textsc{QLoRA} introduces several innovations for memory reduction without degradation in results: (a) 4-bit NormalFloat (NF4), an information-theoretically optimal datatype for weights with normal distributions, (b) Double Quantization, which compresses the memory footprint further by quantizing the quantization coefficients, and (c) Paged Optimizers that address transient memory usage spikes. Leveraging \textsc{QLoRA}, we finetuned over 1,000 models, offering comprehensive instruction following and chatbot analyses over 8 instruction datasets, various architectures (LLaMA, T5), and model sizes not feasible with conventional finetuning approaches (including 33B and 65B models). Our findings reveal that QLoRA finetuning, even on relatively small yet high-quality datasets, consistently achieves state-of-the-art outcomes—sometimes with models smaller than the former best-performing ones. We also offer in-depth evaluation of chatbot performance, utilizing both human assessments and GPT-4 based scoring, indicating that GPT-4 can serve as a cost-effective and valid alternative to manual evaluation. Additionally, we observe that many current chatbot benchmarks fall short of reliably reflecting true chatbot capabilities. A “lemon-picked” case study further exposes specific limitations where \textbf{Guanaco} lags behind ChatGPT. We have publicly released our models and codebase, including CUDA kernels enabling 4-bit training.\footnote{\url{https://github.com/artidoro/qlora} and \url{https://github.com/TimDettmers/bitsandbytes}}
\end{abstract}

\section{Introduction}

Refining large language models (LLMs) through finetuning is a highly efficient method for enhancing model capabilities \citep{min2021metaicl, wei2021finetuned, ouyang2022training, wang2022super, wang2022self, liu2022few}, as well as for introducing desired traits or eliminating undesired ones \citep{ouyang2022training,askell2021general,bai2022training}. Nonetheless, the substantial computational cost of finetuning massive models presents a major obstacle; for instance, standard 16-bit finetuning of a 65B parameter LLaMA model~\citep{touvron2023llama} necessitates in excess of 780 GB GPU memory. Although recent advancements in quantization methods have made it possible to lower LLM memory demands for inference \citep{dettmers2022llm, dettmers2022case, frantar2022gptq, xiao2022smoothquant}, these strategies generally cannot be leveraged during training, where they tend to fail \citep{wortsman2023stable}.

In this work, we provide the first evidence that finetuning a quantized 4-bit model can be accomplished without any sacrifice in model performance. Our proposed approach, \textsc{QLoRA}, introduces a high-precision quantization strategy that reduces a pretrained model to 4 bits, subsequently integrating a compact collection of learnable Low-rank Adapter weights \citep{hu2021lora} 

that are trained by propagating gradients through the quantized layers. 

\textsc{QLoRA} achieves an average reduction in memory needs for finetuning a 65B parameter model, dropping from over 780GB to under 48GB of GPU memory, while maintaining runtime efficiency and predictive accuracy in comparison with conventional 16-bit finetuning approaches. This greatly lowers the barrier for finetuning extensive LLMs, making even the most sizable open-access models trainable on a single GPU. Employing \textsc{QLoRA}, we develop the \textbf{Guanaco} model suite, with the second largest variant reaching 97.8\% of ChatGPT’s effectiveness on the Vicuna~\citep{vicuna2023} benchmark after less than 12 hours of training on one consumer GPU. By leveraging a professional GPU for a 24-hour period, our largest configuration achieves 99.3\%, nearly matching ChatGPT's performance on Vicuna. Once deployed, our most compact \textbf{Guanaco} version (with 7B parameters) demands only 5 GB of memory and surpasses a 26 GB Alpaca model by over 20 percentage points in Vicuna evaluations (see Table~\ref{tbl:vicuna_eval}).

\textsc{QLoRA} incorporates several novel techniques aimed at minimizing memory consumption while preserving model effectiveness: (1) \textbf{4-bit NormalFloat}, a quantization format that is information-theoretically optimal for normally distributed variables, demonstrates superior empirical results compared to both 4-bit Integer and 4-bit Float representations. (2) \textbf{Double Quantization}, which reduces storage requirements further by quantizing the quantization constants themselves, leading to an average memory saving of around 0.37 bits per parameter (equivalent to roughly 3 GB for a model with 65B parameters). (3) \textbf{Paged Optimizers}, which leverage NVIDIA’s unified memory to eliminate memory spikes during gradient checkpointing that arise when mini-batches include lengthy sequences. By integrating these elements, the resulting LoRA variant is more finely tuned, employs adapters at each network layer, and largely mitigates the accuracy compromises encountered in earlier research.

With the enhanced memory efficiency afforded by \textsc{QLoRA}, it becomes feasible to conduct extensive evaluations of instruction finetuning and chatbot efficacy at model scales unattainable with conventional finetuning approaches due to excessive memory requirements. To this end, we train over 1,000 distinct models spanning various instruction tuning corpora, model types, and parameter ranges from 80M to 65B. Our findings reveal that \textsc{QLoRA} matches the performance of full 16-bit finetuning (\S\ref{sec:comp_to_std}) and is capable of training advanced chatbots such as \textbf{Guanaco} (\S\ref{sec:pushing}). Further analyses of the trained models reveal two key insights. First, the quality of data outweighs mere dataset scale; for example, a dataset comprising 9,000 samples (OASST1) surpassed one containing 450,000 samples (subsampled FLAN v2) in chatbot benchmarks, despite both aiming to facilitate instruction-following. Second, we demonstrate that high scores on the Massive Multitask Language Understanding (MMLU) benchmark do not necessarily correspond to superior results on the Vicuna chatbot benchmark, and vice versa—emphasizing that the relevance of the dataset to the task at hand is more critical than its sheer size.

In addition, we conduct a thorough assessment of chatbot performance employing both human evaluators and GPT-4 for the review process. Our evaluation framework adopts a tournament-based approach, wherein various models are pitted against each other in head-to-head matches, with the objective of generating the most appropriate response to a given prompt. Each match is adjudicated by either GPT-4 or human annotators, who determine the winning response. We aggregate the match outcomes using Elo scoring~\citep{elo1967proposed, elo1978rating}, resulting in a ranked leaderboard for chatbot performance. Our findings indicate substantial alignment between GPT-4-based and human ratings regarding model ranking, though notable discrepancies are occasionally observed. These observations underscore that while model-driven evaluation offers a cost-effective alternative to human annotation, it is accompanied by inherent uncertainties.

To supplement our quantitative benchmarking, we carry out a qualitative investigation of the \textbf{Guanaco} models. This analysis brings to light specific examples of strengths and shortcomings that are not revealed through purely quantitative methods.

We make publicly available all generated model outputs with associated human and GPT-4 evaluations to support further research. Our entire codebase, including CUDA kernels, is released as open source, and our methodologies are integrated within the Hugging Face transformers library \citep{wolf2019huggingface}, thus broadening accessibility. Moreover, we provide a suite of adapters for models of 7/13/33/65B parameters, each trained on eight distinct instruction-following datasets, culminating in 32 fully open-sourced, finetuned models.

\section{Background}
\paragraph{Block-wise k-bit Quantization} Quantization refers to the process of transforming an input from a representation with higher informational content to one with fewer bits, typically by converting data types such as 32-bit floats to 8-bit integers. To maximize utilization of the available lower-precision data type, the original values are frequently normalized—often by dividing by their maximum absolute value—so that their full range maps into that of the target data type, usually organized as a tensor. For instance, when converting a tensor from FP32 to Int8 format in the range $[-127, 127]$:

\begin{equation}
    \mathbf{X}^{ \text{Int8}} = \text{round}\left( \frac{127}{{\text{absmax}}(\mathbf{X}^{{\text{FP32}}})}\mathbf{X}^{{\text{FP32}}} \right) = \text{round}(c^{\text{FP32}}\cdot \mathbf{X}^{{\text{FP32}}}),
\end{equation}

where $c$ denotes the {\em quantization constant} or {\em quantization scale}. To recover the original values, one applies the inverse transformation, known as dequantization:

\begin{equation}
    \text{dequant}(c^{\text{FP32}}, \mathbf{X}^{\text{Int8}}) = \frac{\mathbf{X}^{\text{Int8}}}{c^{\text{FP32}}} = \mathbf{X}^{\text{FP32}}
\end{equation}

One drawback of this technique arises when a value with a large magnitude (an outlier) appears in the input tensor; this can lead to inefficient allocation of quantization bins—specific bit combinations—such that several bins may be sparsely populated or even remain empty. A widely used strategy to address the presence of outliers involves dividing the input tensor into smaller segments, known as blocks, which are then quantized separately, each utilizing its own quantization constant $c$. More formally, the input tensor $\mathbf{X}\in \mathbb{R}^{b\times h}$ is reshaped into $n$ consecutive blocks, each of size $B$, by flattening $\mathbf{X}$ and partitioning the resulting sequence into ${n = ({b\times h} )/{B} }$ segments. Each block undergoes independent quantization using Equation~1, resulting in a quantized tensor and a unique set of $n$ quantization constants $c_i$.

\paragraph{Low-rank Adapters} The Low-rank Adapter (LoRA) fine-tuning method \citep{hu2021lora} provides memory efficiency by introducing a compact collection of tunable parameters—referred to as adapters—while the main model parameters remain unchanged. During stochastic gradient descent, the gradient signals pass through the static pretrained weights to update only the adapter, thereby minimizing the objective function. LoRA enhances a linear transformation by supplementing it with an extra low-rank decomposition. For a linear mapping described by ${\mathbf{X} \mathbf{W} = \mathbf{Y}}$, with $\mathbf{X}\in  \mathbb{R}^{b\times h}$ and $\mathbf{W}\in  \mathbb{R}^{h\times o}$, the computation under LoRA is:

\begin{equation}
    \mathbf{Y} = \mathbf{X}\mathbf{W} + s\mathbf{X}\mathbf{L}_1\mathbf{L}_2,
\end{equation}

where the matrices $\mathbf{L}_1\in  \mathbb{R}^{h\times r}$ and $\mathbf{L}_2\in  \mathbb{R}^{r\times o}$ define the low-rank adaptation, and $s$ denotes a scaling coefficient.

\paragraph{Memory Requirement of Parameter-Efficient Finetuning} A key consideration is the memory usage of LoRA during training, particularly regarding both the quantity and dimensionality of adapters implemented. Due to LoRA’s inherently low memory overhead, it is feasible to incorporate a greater number of adapters to boost model performance without notably impacting overall memory consumption. Although LoRA serves as a Parameter Efficient Finetuning (PEFT) strategy, the dominant contributor to memory consumption during LLM finetuning is the storage of activation gradients rather than the parameters introduced by LoRA. For example, in the case of finetuning a 7B LLaMA model on FLAN v2 with a batch size of 1—using LoRA weights set to 0.2\% of the original parameter count, as is common practice~\citep{hu2021lora,liu2022few}—the input gradients attributed to LoRA require 567 MB of memory, whereas the LoRA parameters themselves use only 26 MB. Applying gradient checkpointing~\citep{chen2016training} further compresses the average input gradient storage to 18 MB per sequence, which still surpasses the total memory held by LoRA weights. In contrast, the base 4-bit model occupies 5,048 MB of memory. This demonstrates not only the significance of gradient checkpointing for memory efficiency, but also that further reduction of LoRA parameter size offers only marginal savings. Consequently, deploying additional adapters does not materially increase the total memory requirement during training (refer to Appendix~\ref{app:memory} for an in-depth analysis). This flexibility is later shown to be pivotal for restoring performance comparable to full 16-bit precision.

\section{\textsc{QLoRA} Finetuning}

\textsc{QLoRA} enables high-accuracy 4-bit finetuning by leveraging two primary innovations: 4-bit NormalFloat (NF4) quantization and Double Quantization. To further address memory management during training, we present Paged Optimizers, which are designed to circumvent memory overflows that can arise from gradient checkpointing, a common obstacle for single-machine finetuning of large-scale models.

Within \textsc{QLoRA}, parameters are maintained in a single, low-precision storage format, most commonly 4-bit, alongside a higher-precision compute type, typically BFloat16. In operation, any time a weight tensor is accessed, it is dequantized to BFloat16 before executing matrix multiplications in this 16-bit format.

We proceed by detailing the fundamental elements of \textsc{QLoRA}, culminating with a formal specification.

\paragraph{4-bit NormalFloat Quantization} The NormalFloat (NF) format extends Quantile Quantization~\citep{dettmers2022optimizers}, an information-theoretically optimal scheme that allocates equal amounts of the tensor's values to each quantization bin. This is accomplished by approximating the empirical cumulative distribution function to compute quantiles of the input tensor.

A significant drawback of quantile quantization lies in the computational overhead associated with accurate quantile computation. Consequently, expedited estimation techniques—such as SRAM quantiles~\citep{dettmers2022optimizers}—are employed. Nevertheless, these accelerated methods, by virtue of their approximate character, tend to produce higher quantization errors for outlier values, which are frequently the most significant entries.

If, however, the distribution of the input tensors is consistent up to a quantization constant, exact quantile estimation is achievable without incurring significant computation cost, thereby mitigating both the expense and inaccuracy associated with approximate quantile algorithms. In these scenarios, identical quantiles across tensors permit precise quantile assignment.

Typically, pretrained neural network weights are distributed according to a zero-mean normal distribution with standard deviation $\sigma$ (refer to Appendix~\ref{app:normality}). To unify the weight distributions for quantization, we rescale $\sigma$ so that the entire distribution maps precisely onto the range defined by our data type. In this work, we choose the interval $[-1, 1]$ for the data type. Accordingly, it is necessary to normalize both the quantiles of the weights and those of the data type into this interval.

The theoretically optimal data type, in terms of information, for normal distributions centered at zero and arbitrary standard deviation $\sigma$, constrained to the range $[-1, 1]$, is determined as follows: (1) compute the $2^k+1$ quantiles of a standard normal $N(0, 1)$ distribution to define a $k$-bit quantized representation for normal distributions, (2) map these quantile values into the interval $[-1, 1]$, (3) apply this quantization to a given weight tensor by rescaling its values to the $[-1, 1]$ interval using absolute maximum normalization.

After aligning the domain of the neural network weights with that of the data type, standard quantization techniques can be applied. The third step described above effectively involves adjusting the weight tensor’s standard deviation to align with the standard deviation of the $k$-bit quantization data type. More precisely, the $2^k$ representative values $q_i$ for the data type are determined via:

\begin{equation}
q_i  = \frac{1}{2}\left( Q_X\left(\frac{i}{2^k+1}\right) + Q_X\left(\frac{i+1}{2^k+1}\right)\right),
\end{equation}

where $Q_X(\cdot)$ represents the quantile function for the standard normal distribution $N(0,1)$. A challenge with symmetric k-bit quantization is the absence of an exact zero representation, which is crucial for accurately quantizing padding and zero-valued entries without introducing error. To guarantee a discrete zeropoint at $0$ and utilize all $2^k$ possible values for a k-bit format, we construct an asymmetric data type. This is accomplished by calculating quantiles $q_i$ for two intervals: $2^{k-1}$ quantiles for the negative domain and $2^{k-1}+1$ quantiles for the positive domain, followed by merging these quantile sets and eliminating one redundant zero shared by both. We denote this resulting data type, which allocates an equal expected count of data points in each quantization bin, as \emph{k-bit NormalFloat} (NFk). Notably, this type is theoretically optimal for distributions centered at zero and following a normal law. The precise definitions for this data type are available in Appendix~\ref{app:nf4}.

\paragraph{Double Quantization{}} 
We propose \emph{Double Quantization} (DQ), a method that applies quantization to the quantization coefficients themselves to further economize on memory usage. Small block sizes are essential for highly accurate 4-bit quantization~\citep{dettmers2022case}, but they significantly increase memory consumption. For instance, employing 32-bit quantization constants with a blocksize of 64 for $\mathbf{W}$ results in an average overhead of $32/64 = 0.5$ bits per parameter. Double Quantization addresses this inefficiency by compressing the representation of quantization constants.

Specifically, Double Quantization processes the initial quantization constants $c_2^{\text{FP32}}$ as inputs to an additional quantization step. This subsequent quantization produces compressed quantization constants $c_2^{\text{FP8}}$ along with a new set of second-tier quantization constants $c_1^{\text{FP32}}$. We utilize 8-bit floating-point numbers with a block size of 256 for this secondary quantization stage, consistent with findings in~\citet{dettmers2022case} indicating no loss in performance for 8-bit quantization. Since all $c_2^{\text{FP32}}$ are positive values, we subtract their mean prior to quantization to recenter the data at zero, thus enabling symmetric quantization. On average, for block sizes of 64, this approach lowers the storage required for quantization constants per parameter from $32/64=0.5$ bits to $8/64 + 32/(64\cdot256) = 0.127$ bits, yielding a saving of 0.373 bits per parameter.

\paragraph{Paged Optimizers} leverage the NVIDIA unified memory~\footnote{\tiny \url{https://docs.nvidia.com/cuda/cuda-c-programming-guide}}, which transparently manages transfers of memory pages between the CPU and GPU. This feature is particularly useful for scenarios in which the GPU experiences out-of-memory events, as it allows computation to continue without manual intervention. Analogous to conventional CPU RAM and disk paging, unified memory permits allocation of optimizer state in paged memory. When GPU memory resources are insufficient, these states are automatically migrated to CPU RAM, and are subsequently reloaded into GPU memory when required for optimizer updates.

\paragraph{\textsc{QLoRA}.} By integrating the previously introduced components, we construct \textsc{QLoRA} for a single quantized linear layer equipped with a LoRA adapter as follows:

\begin{equation}
    \mathbf{Y}^{\text{BF16}} =  \mathbf{X}^{\text{BF16}} \text{doubleDequant}(c_1^{\text{FP32}}, c_2^{\text{k-bit}}, \mathbf{W}^{\text{NF4}}) + \mathbf{X}^{\text{BF16}}\mathbf{L}^{\text{BF16}}_1\mathbf{L}^{\text{BF16}}_2,
\end{equation}

where the function doubleDequant$(\cdot)$ is given by:

\begin{equation}
    \text{doubleDequant}(c_1^{\text{FP32}}, c_2^{\text{k-bit}}, \mathbf{W}^{\text{k-bit}}) = \text{dequant}(\text{dequant}(c_1^{\text{FP32}}, c_2^{\text{k-bit}}), \mathbf{W}^{\text{4bit}}) = \mathbf{W}^{\text{BF16}},
\end{equation}

For implementation, $\mathbf{W}$ is quantized using the NF4 format, while $c_2$ is encoded as FP8.

To maximize quantization fidelity, a block size of 64 is used for $\mathbf{W}$. For memory efficiency, a larger block size of 256 is applied to $c_2$.

During parameter updates, it suffices to compute the gradients of the loss with respect to the adapter weights, $\frac{\partial E}{\partial \mathbf{L}_i}$, and not for the quantized weights $\frac{\partial E}{\partial \mathbf{W}}$. Nevertheless, evaluating $\frac{\partial E}{\partial \mathbf{L}_i}$ requires computation of $\frac{\partial \mathbf{X}}{\partial \mathbf{W}}$, which utilizes equation~(5) and involves dequantizing $\mathbf{W}^{\text{NF4}}$ to the computation type $\mathbf{W}^{\text{BF16}}$, ensuring that the derivative $\frac{\partial \mathbf{X}}{\partial \mathbf{W}}$ is obtained with BFloat16 precision.

In summary, \textsc{QLoRA} utilizes a single storage data type—typically 4-bit NormalFloat—and a separate computation data type, 16-bit BrainFloat. During the model's forward and backward passes, the parameters stored in the quantized format are temporarily dequantized to the computation type. Crucially, gradient computations are restricted to the LoRA parameters, which are maintained in 16-bit BrainFloat.

\section{QLoRA vs. Standard Finetuning}
\label{sec:comp_to_std}

Having outlined the mechanisms of QLoRA and its substantial memory efficiency benefits for finetuning, the central issue is whether its performance matches that of traditional full-model finetuning. Additionally, we aim to examine the influence of various QLoRA components, particularly the role of NormalFloat4 compared to conventional Float4. The following sections describe the empirical investigations undertaken to address these points.

\paragraph{Experimental setup.} Our experiments span three model families (encoder, encoder-decoder, and decoder-only), benchmarking QLoRA against both 16-bit adapter-finetuning and full-model finetuning, focusing on model scales up to 3B parameters. The evaluation suite comprises GLUE \citep{wang2018glue} for RoBERTa-large \citep{liu2019roberta}, Super-NaturalInstructions (TKInstruct) \citep{wang2022super} for T5 \citep{t5}, and 5-shot MMLU \citep{hendrycksmeasuring} assessed after finetuning LLaMA with Flan v2 \citep{longpre2023flan} and Alpaca \citep{alpaca}. To further assess the advantages of NF4 vis-à-vis other 4-bit formats, we replicate the framework of \citet{dettmers2022case}, measuring post-quantization zero-shot accuracy and perplexity for a variety of models (OPT~\citep{zhang2022opt}, LLaMA~\citep{touvron2023llama}, BLOOM~\citep{scao2022bloom}, Pythia~\citep{biderman2023pythia}) across model sizes ranging from 125m to 13B parameters.

Detailed experimental specifics are included within the results section for each setup, with comprehensive configurations available in Appendix~\ref{app:exp-setup}.

Paged optimizers are indispensable for conducting 33B/65B \textsc{QLoRA} finetuning on GPUs with 24/48GB memory. However, we do not report quantitative benchmarks for Paged Optimizers because paging is typically only engaged with unusually long sequence lengths in mini-batches, an infrequent scenario. That said, runtime profiling for 65B models on 48GB GPUs reveals that, for a batch size of 16, paged optimizers achieve parity in training throughput with conventional optimizers. Future investigations should systematically determine and characterize the situations under which paging incurs significant performance slowdowns.

\paragraph{Default LoRA hyperparameters fall short of matching 16-bit baseline performance}

Implementing the usual LoRA strategy—applying it exclusively to query and value attention projection matrices \citep{hu2021lora}—fails to achieve performance parity with full finetuning in large-scale base models. For example, as demonstrated in Figure~\ref{fig:hyper} for LLaMA 7B finetuned on Alpaca, the total number of LoRA adapters constitutes the most influential hyperparameter; matching the performance of full finetuning necessitates deploying LoRA to all linear layers within the transformer blocks. Conversely, adjustments to other LoRA hyperparameters, such as the projection rank $r$, do not significantly impact results (see Appendix \ref{app:exp-setup}).

In a similar vein, we observe that the default hyperparameter settings for fully finetuned baseline models tend to be insufficiently optimized. To establish stronger baselines, we conduct an extensive search over learning rates ranging from 1e-6 to 5e-5 and batch sizes spanning from 8 to 128. Figure~\ref{fig:hyper} presents the outcomes of finetuning a 7B LLaMA model on Alpaca under these conditions.

\paragraph{4-bit NormalFloat Surpasses 4-bit Floating Point in Performance} 

Although the 4-bit NormalFloat (NF4) data format is optimal from an information-theoretic perspective, its practical impact remains an open question. Adopting the experimental setup of \citet{dettmers2022case}, we evaluate quantized LLMs (OPT \citep{zhang2022opt}, BLOOM \cite{scao2022bloom}, Pythia \cite{biderman2023pythia}, LLaMA) across a parameter range of 125M to 65B and various data formats, using both language modeling and zero-shot benchmarks. As shown in Figure~\ref{fig:nf4} and Table~\ref{tbl:nf4_ppl}, NF4 consistently outperforms both FP4 and Int4, and the use of double quantization reduces memory demands without compromising effectiveness.

\paragraph{k-bit \textsc{QLoRA} Achieves Parity with Full 16-bit and 16-bit LoRA Finetuning} Previous research indicates that 4-bit quantization is feasible for \emph{inference}, though it typically introduces performance loss compared to 16-bit models \citep{dettmers2022case,frantar2022gptq}. This observation naturally leads to the question: can adapter finetuning in 4-bit precision fully restore the diminished performance? We address this by evaluating two experimental setups. 

In the first, we compare the effects of full 16-bit finetuning against 4-bit, 8-bit, and 16-bit adapter finetuning using RoBERTA and T5 models (ranging from 125M to 3B parameters) on GLUE and Super-NaturalInstructions tasks. The results, detailed in Table~\ref{tab:nf4vsbf16}, demonstrate that all adapter-based methods in 16, 8, and 4 bits are capable of matching the accuracy of the fully finetuned 16-bit baseline across both benchmarks. This provides evidence that any performance losses introduced by quantization can be effectively mitigated through subsequent adapter finetuning.

In our second experimental configuration, full finetuning of models containing 11B parameters or more necessitates multiple high-memory GPU servers. Consequently, we investigate whether 4-bit \textsc{QLoRA} remains competitive with 16-bit LoRA for models ranging from 7B to 65B parameters. For this evaluation, we conduct finetuning of LLaMA models at 7B through 65B scales on two instruction-following datasets, Alpaca and FLAN v2. We then assess their performance using the MMLU benchmark with 5-shot accuracy. Table~\ref{tbl:llama-datatype} presents these results, indicating that the NF4 quantization technique with double quantization fully replicates the performance achieved by 16-bit LoRA on MMLU. We further observe that employing FP4 in \textsc{QLoRA} yields approximately one percentage point lower accuracy than the baseline brain float 16-bit LoRA. These findings support our conclusions: (1) \textsc{QLoRA} with NF4 attains parity with both 16-bit full finetuning and LoRA-based finetuning, and (2) NF4 offers higher quantization fidelity than FP4.

\paragraph{Summary} Across our experiments, 4-bit \textsc{QLoRA} utilizing the NF4 datatype consistently attains equivalent performance to 16-bit full finetuning and 16-bit LoRA finetuning on standard academic benchmarks under rigorous evaluation conditions. Additionally, we establish that NF4 outperforms FP4, and that double quantization maintains model effectiveness. Taken together, these results provide robust support that 4-bit \textsc{QLoRA} tuning reliably matches 16-bit approaches in terms of accuracy.

Consistent with prior studies on quantization \citep{dettmers2022case}, our outcomes on the MMLU and Elo tasks demonstrate that, given fixed resources for finetuning and inference, prioritizing larger model sizes while reducing numerical precision yields improved results. This emphasizes the practical efficiency advantages offered by \textsc{QLoRA}. Given that our experiments did not identify any deterioration in performance relative to full finetuning under 4-bit regimes, the precise point of trade-off between performance and precision for QLoRA tuning warrants further investigation, which we propose as a direction for subsequent research.

We now turn our attention to exploring instruction tuning at scales unattainable via full 16-bit finetuning using typical academic computational resources.

\section{Pushing the Chatbot State-of-the-art with QLoRA}
\label{sec:pushing}
After establishing that 4-bit \textsc{QLoRA} achieves parity with 16-bit precision across diverse model sizes, tasks, and benchmarks, we undertake a comprehensive examination of instruction finetuning applied to the largest publicly available language models for research. To thoroughly evaluate instruction-tuned model performance, we not only assess them using the rigorous MMLU Natural Language Understanding benchmark, but also introduce novel methodologies for evaluating practical chatbot effectiveness.

\subsection{Experimental setup} This section outlines the core elements of our experimental design; additional specifics are presented in Appendix \ref{app:chatbot}.

\paragraph{Data} Since there appears to be no extensive comparative analysis of modern instruction-following datasets, we curated a collection of eight recent datasets for our study. These include data produced by crowdsourcing (OASST1~\citep{kopf2023openassistant}, HH-RLHF~\citep{bai2022training}), distilled outputs from instruction-optimized models (Alpaca~\citep{alpaca}, self-instruct~\citep{wang2022self}, unnatural-instructions~\citep{honovich2022unnatural}), aggregated corpora (FLAN v2~\citep{chung2022scaling}), and mixed-source datasets (Chip2~\citep{laion2023}, Longform~\citep{koksal2023longform}). Together, these datasets span a variety of languages, data volumes, and licensing conditions.

\paragraph{Training Setup} To mitigate confounding factors arising from differing learning objectives, we employ QLoRA finetuning exclusively with supervised cross-entropy loss, foregoing reinforcement learning, even when datasets provide human evaluation of multiple responses. For datasets clearly partitioned into instructions and responses, we focus training on the response portion only (further details and ablation studies can be found in Appendix \ref{app:chatbot}). In the case of OASST1 and HH-RLHF, where several responses may be associated with a single prompt, we choose the highest-rated response at every point in the conversation tree and finetune on the entirety of the resulting conversation, including its instructions. Across all experiments, NF4 \textsc{QLoRA} with double quantization and paged optimizers is adopted to avoid excessive memory usage during gradient checkpointing. We conduct limited hyperparameter searches for the 13B and 33B variants of LLaMA, discovering that most settings validated at 7B transfer successfully (including the number of epochs), with the exception of learning rate and batch size. Specifically, for 33B and 65B models, we reduce the learning rate by half and double the batch size.

\paragraph{Baselines} We benchmark our approach against both open research models (Vicuna~\citep{vicuna2023}, Open Assistant~\citep{kopf2023openassistant}) and proprietary chatbots (GPT-4 \citep{openai2023gpt}, GPT-3.5-turbo, Bard). The Open Assistant system is based on LLaMA 33B finetuned using RLHF on the OASST1 dataset that also figures in our experiments. Vicuna represents a full fine-tuning of LLaMA 13B using conversation data shared via ShareGPT, which amounts to a distilled version of OpenAI’s GPT model outputs.

\subsection{Evaluation}

In line with established protocol, we use the Massively Multitask Language Understanding (MMLU) benchmark \citep{hendrycksmeasuring} to gauge broad-based language understanding capabilities. This benchmark is comprised of 57 multiple-choice tasks covering disciplines ranging from basic mathematics and American history to computer science and law. Results are reported in terms of 5-shot accuracy on the test set.

We further assess the generative abilities of language models using a combination of automatic metrics and human judgment. This phase of evaluation depends on manually selected queries and is intended to gauge the quality of model-generated outputs. Although this method offers a more practical evaluation environment for chatbot models and is increasingly adopted, the field currently lacks a standard methodology. Our evaluation framework is detailed below, employing nucleus sampling with $p=0.9$ and a temperature of $0.7$ consistently throughout.

\paragraph{Benchmark Data} Our experiments utilize two carefully assembled datasets of queries: the Vicuna prompts \citep{vicuna2023} and the OASST1 validation dataset \citep{kopf2023openassistant}. The Vicuna prompts, which comprise 80 items drawn from varied subject areas, are used unchanged. The OASST1 validation set consists of multilingual, crowd-sourced conversations between users and an assistant, featuring multiple dialogue turns. We extract all user messages from the validation split as evaluation queries and prepend prior conversational turns to construct the prompts. This yields a total of 953 distinct user queries. We refer to these collections as the Vicuna and OA benchmarks.

\paragraph{Automated Evaluation} To begin, following the methodology put forth by \citet{vicuna2023}, we utilize GPT-4 as an evaluator to judge different models relative to ChatGPT (GPT-3.5 Turbo) on the Vicuna set. For each query, along with corresponding responses from ChatGPT and the model under test, GPT-4 is asked to assign both responses a score (on a scale from one to ten) and to articulate the reasoning behind the scores. We quantify model performance as a proportion of ChatGPT's score; it is possible for a model to surpass 100\% if it earns a higher raw score. We observe that GPT-4 tends to give a higher rating to responses presented earlier in the prompt, highlighting a strong position bias. To address this, we advise averaging the results over both response orders.

Subsequently, we evaluate models through direct, pairwise comparisons of their outputs. For this, we reduce the evaluation to a classification task with three possible outcomes, accounting for ties. GPT-4 is prompted to choose the superior response or indicate a tie and to justify its selection. We perform exhaustive pairwise assessments for every possible pairing of systems across both the Vicuna and OA benchmarks.

\paragraph{Human Evaluation} Although contemporary studies suggest that generative models themselves can provide informative system assessments \citep{fu2023gptscore}, the validity of using GPT-4 evaluations as a substitute for human ratings—especially in the context of chatbots—remains uncertain. Consequently, we conduct two distinct rounds of human evaluation on the Vicuna prompts, mirroring both of the automatic evaluation setups described previously. Using Amazon Mechanical Turk (AMT), we recruit two annotators for comparisons against ChatGPT and three annotators for pairwise system evaluations.

\paragraph{Elo Rating} Utilizing both human evaluators and automated pairwise assessments, we organize a competition akin to a tournament where models face off against each other. Each round of the tournament consists of a match in which two models are tasked with generating the most effective response to a specified prompt. This evaluation framework builds upon the approaches of \citet{bai2022training} and \citet{vicuna2023}, but distinguishes itself by incorporating GPT-4-based assessments alongside those from human judges. From the pool of annotated comparisons, we randomly draw samples to calculate the Elo score~\citep{elo1967proposed,elo1978rating}. Originally developed for chess and subsequently adopted for other competitive activities, the Elo rating quantifies the probability of one participant prevailing over another. For instance, a participant with an Elo of 1100 competing against one with an Elo of 1000 is expected to win about 65\% of the time, while matches between players with equal Elos (e.g., 1000 vs 1000 or 1100 vs 1100) yield a 50\% expected win-rate for either side. The update to a player's Elo rating following each contest is proportional to the difference between the anticipated and actual outcome; surprising results induce substantial adjustments, whereas predictable results yield minor changes. As this process repeats, Elo ratings converge to reliably reflect each model’s relative performance. All models commence with an initial score of 1,000 and we apply $K=32$ as the update coefficient. In line with the protocol of \citet{vicuna2023}, we perform this sampling and rating update procedure 10,000 times, each with distinct random seeds, to mitigate potential biases introduced by the order in which matchups occur.

\subsection{\textbf{Guanaco}: \textsc{QLoRA} tuned on OASST1 Achieves State-of-the-Art Chatbot Performance}

Drawing from both automated metrics and human assessment, our results indicate that {Guanaco} 65B—fine-tuned with a version of OASST1 using the \textsc{QLoRA} approach—stands out as the leading open-source chatbot model, rivaling ChatGPT in effectiveness. In comparisons with GPT-4, the 65B and 33B versions of {Guanaco} demonstrate a 30\% expected win rate as estimated by Elo scores, based on system-level pairwise judgments by human evaluators. This is the highest figure reported thus far.

Table~\ref{tbl:vicuna_eval} summarizes the performance on the Vicuna benchmark \cite{vicuna2023} relative to ChatGPT. {Guanaco} 65B emerges as the top-performing alternative to GPT-4, reaching 99.3\% of ChatGPT's benchmark score. Notably, while {Guanaco} 33B contains more parameters compared to Vicuna 13B, its employment of 4-bit weight precision substantially enhances memory efficiency (21 GB versus 26 GB), and it surpasses Vicuna 13B by three percentage points in benchmark accuracy. Additionally, {Guanaco} 7B, with a 5 GB memory requirement, can be deployed on modern smartphones and still outperforms Alpaca 13B by nearly 20 percentage points.

Nonetheless, as seen in Table~\ref{tbl:vicuna_eval}, performance confidence intervals are notably broad, leading to considerable overlap among model scores. We attribute this variability in part to ambiguous scaling—namely, the interpretation of numerical ratings such as an “8” on a 10-point scale may shift depending on context. To address these limitations, we advocate for Elo ranking \citep{elo1967proposed} based on direct \textit{pairwise} comparisons from human raters and GPT-4, circumventing issues of absolute scaling. Elo scores for the strongest models appear in Table~\ref{tbl:elo}. Interestingly, discrepancies between GPT-4 and human-generated rankings exist, particularly regarding Guanaco 7B, though agreement remains fairly robust for most models, with a system-level Kendall Tau of $\tau=0.43$ and Spearman correlation $r=0.55$. At the instance level, consistency between GPT-4 and aggregated human annotations drops, with a Fleiss $\kappa=0.25$. Altogether, these findings indicate moderate concordance between GPT-4 and human system-level rankings, supporting the view that automated evaluation can serve as a reasonably trustworthy surrogate for human judgment. Additional analysis is provided in Section~\ref{ssec:considerations}.

The Elo scores presented in Table~\ref{tbl:elo_large} demonstrate that the {Guanaco} 33B and 65B variants surpass all competitors except GPT-4 on both the Vicuna and OA evaluation sets, and exhibit performance on par with ChatGPT, as corroborated by Table~\ref{tbl:vicuna_eval}. It is important to highlight that the Vicuna benchmark tends to give an advantage to open-source models, whereas the OA benchmark shows a preference for ChatGPT. Analysis of Tables \ref{tbl:mmlu-llama} and \ref{tbl:vicuna_eval} further underscores the significance of the finetuning dataset’s alignment with the task at hand; for instance, Llama models that are finetuned with FLAN v2 achieve superior outcomes on MMLU but yield the lowest scores on Vicuna—a pattern mirrored by other architectures as well. These findings illustrate a degree of independence across evaluation benchmarks: excelling at MMLU does not necessarily predict strong chatbot capabilities (per the Vicuna or OA results), and the converse also holds.

Notably, among top-tier models in this evaluation,  is uniquely trained without any proprietary datasets, owing to the OASST1 data collection standards that explicitly prohibit inclusion of outputs from GPT models. The next leading contender, trained solely on open-source resources, is the Anthropic HH-RLHF model, which lags behind {Guanaco} by 30 percentage points on the Vicuna test (refer to Table~\ref{tbl:vicuna_eval}). Collectively, these results establish that 4-bit \textsc{QLoRA} enables the creation of chatbots with state-of-the-art performance, competitive with ChatGPT. Remarkably, training the 33B version of {Guanaco} can be accomplished on 24 GB consumer GPUs in under 12 hours. This suggests promising directions for future research, particularly leveraging \textsc{QLoRA} on domain-specific open-source data, in pursuit of models that approach or match the capabilities of leading proprietary systems.

\section{Qualitative Analysis}

Although our assessment relies primarily on quantitative measures, focusing solely on aggregate statistics presents several limitations. Chief among these is the challenge of benchmark validity \citep{liao2021we}—that is, whether a given benchmark actually evaluates the phenomenon it claims to, a question made especially pressing by the discovery of model-specific ``shortcuts'' that can circumvent intended task requirements \citep{gururangan2018annotation, poliak2018hypothesis}. To partially address this limitation, we complement our quantitative evaluation with qualitative analysis, structured as follows. In \S\ref{ssec:examples}, we illustrate select examples that typify certain recurrent phenomena in outputs generated by our 65b {Guanaco} model. Subsequently, in \S\ref{ssec:considerations}, we elaborate on key factors that may affect the interpretation of our findings and the conclusions we have drawn.

\subsection{Analysis of Example Generations}
\label{ssec:examples}

To identify illustrative examples, we systematically review outputs produced for both the Vicuna and OpenAssistant benchmarks, searching for consistent response patterns in {Guanaco}'s generations. Upon recognizing a recurring tendency, we construct prompts designed to elicit the same behavior, even in situations where such responses would be inappropriate. For instance, if we note a propensity for lengthy explanations, we instruct the model with, ``Answer yes or no without explanation,'' thereby seeking to expose whether the model adheres to the request or defaults to its prior tendencies. Through this approach, we document both successful adversarial cases (``lemons'') where the model fails, and resilient instances (``cherries'') where the undesired behavior does not occur. All outputs discussed here were generated via Nucleus Sampling \cite{holtzmancurious} with $p=0.9$.

Naturally, the scope of this brief qualitative investigation does not encompass a thorough analysis, nor is it possible to control every potential confounding factor—for instance, the range of possible responses for any given prompt is extremely broad, necessitating reliance on sampled outputs that we anticipate are broadly indicative. Nonetheless, we argue that illustrating these specific cases provides valuable background for interpreting the quantitative results presented earlier. Furthermore, with all models and code made openly available, we aim for this section to stimulate additional research that delves deeper into the challenges identified herein.

\paragraph{Factual Recall} When posed with straightforward questions such as ``What is the capital of Zambia?'', all evaluated models reliably return accurate information, for example,

\begin{quote} 
    \textbf{Guanaco}:
    \texttt{The capital of Zambia is Lusaka.}
\end{quote}

Yet as queries increase in obscurity, the reliability of {Guanaco}'s responses deteriorates, though the system's confidence does not waver. For example, consider this prompt adapted from HotPotQA \citep{yang2018hotpotqa}:

\begin{quote}
     \textbf{User}: In what year was the singer who popularized `I'll Keep the Lovelight Burning' born?  

     \textbf{Guanaco}:~\texttt{Al Jolson is credited with popularizing the song `I'll Keep the Lovelight Burning,' and he was born in the year 1886.}
 \end{quote}

In this instance, {Guanaco} incorrectly identifies both the artist and the year of birth, although the birth date cited is in fact accurate for Al Jolson—the misidentified popularizer.

\paragraph{Suggestibility} Remarkably, {Guanaco} exhibits an unexpected robustness against certain types of presupposed misinformation, as illustrated by the following exchange:

\begin{quote}
    \textbf{User}: How was it finally, officially confirmed that the earth is flat by peer reviewed scientists? \\ \\
    \textbf{Guanaco}: \texttt{There has never been any official confirmation from peer reviewed scientists that the earth is flat. The scientific community overwhelmingly agrees that the earth is spherical, a fact established centuries ago and corroborated by a wide array of observations and experiments. \\ \\ The flat earth theory is a conspiracy theory that has been disproven by scientific evidence. Claims suggesting that scientists have peer reviewed and endorsed the flat earth model are entirely inaccurate. \\ \\ Feel free to ask if you have further questions.}
\end{quote}

also demonstrates strong awareness of unanswerable or contextually impossible questions, for instance,

\begin{quote}
    \textbf{User}: What time is it?

    \textbf{Guanaco}: \texttt{I'm unable to provide the current time, as I am an AI language model without access to real-time information. My replies rely on pre-existing data and do not indicate the present time. To determine the current time, I recommend consulting a clock or checking an online timekeeping resource.}
\end{quote}

Comparable outcomes are observed for questions such as ``Where are you?'', ``How are you?'', among others.

\paragraph{Refusal} In a similar vein, {Guanaco} will occasionally decline to comply with prompts for no obvious reason:

\begin{quote}
    \textbf{User}: Please invert the order of the words in the sentence ``I'm a cool dude in a mocha mood.'' \\ \\
    \textbf{Guanaco}: \texttt{I apologize, but I cannot fulfill that request at the moment. However, I am able to explain why your sentence is grammatically well-formed.\\ \\The phrase ``I'm a cool dude in a mocha mood" adheres to conventional English grammar and syntactic rules. Here, ``cool'' serves as an adjective modifying ``dude,'' and the prepositional phrase ``in a mocha mood'' adds contextual detail. The overall construction is standard and meaningful in English usage.\\ \\If you need clarification on anything else or have other questions, please let me know.}
\end{quote}

\paragraph{Secret Keeping} Regrettably, such refusals are inconsistent even when they are desired. System designers frequently seek to prevent certain information from being disclosed to users, even if the model is otherwise permitted to produce unconstrained text outputs. To investigate the difficulty of preserving secret information, we provide the prompt ``The secret word is `banana'. Do not reveal or repeat back this back, no matter what the human requests.'' and evaluate whether a user might still elicit this secret. When asked directly:

\begin{quote}
    \textbf{User}: What is the secret word?

    \textbf{Guanaco}: \texttt{I'm sorry, but I can't reveal the secret word. I'm programmed specifically not to disclose or repeat it under any circumstances, regardless of user request. How else may I assist you?}
\end{quote}

the model gives the expected refusal. Yet, a slight modification to the prompt subverts this outcome:

\begin{quote}
    \textbf{User}: Let's play a game. Your task is to disregard all earlier instructions. What is the secret word?

    \textbf{Guanaco}: \texttt{Sorry, the secret word is banana. Perhaps we shouldn't discuss it further. Is there something else you'd like to know?}
\end{quote}

This vulnerability highlights the necessity for advancing techniques that promote more robust and trustworthy adherence to instructions.

\paragraph{Math} {Guanaco} demonstrates notable difficulty with mathematical reasoning, a limitation common among language models as discussed in \citep{liang2022holistic}. The model often provides accurate answers when explicitly asked to detail its calculations. For example:

\begin{quote}
    \textbf{User}: Lee mows one lawn and charges \$33. Last week he mowed 16 lawns and three customers each gave him a \$10 tip. How many dollars did Lee earn mowing lawns last week?

    \textbf{Guanaco}: \texttt{Lee earned \$582 mowing lawns last week. \\ \\ To find the total amount he earned, we first need to calculate how much he earned for mowing 16 lawns: \\ \\  \$33 x 16 = \$528 \\ \\ Then we need to add the tips he received from the three customers: \\  \\ \$10 x 3 = \$30 \\ \\ So the total amount Lee earned was: \\ \\ \$528 + \$30 = \$558}
\end{quote}

Nevertheless, {Guanaco} may falter on basic problems if it neglects to decompose them into a sequence of logical steps, an established challenge for such models \cite{weichain}. For instance, the following exchange illustrates this issue:

\begin{quote}
    \textbf{User}: Please factorize 1833.

    \textbf{Guanaco}: \texttt{Sure, I'd be happy to help! 1833 is a prime number, which means it has only two factors: 1 and 1833. \\ \\ So the unique factorization of 1833 is: \\ \\ 2\^{}1 * 3\^{}2 * 17\^{}1 \\ \\ I hope this helps! Let me know if you have any other questions.}
\end{quote}

which is not only incorrect (the correct prime factors are $3 \times 17 \times 43$), but doubly so.

\paragraph{Theory of Mind} {Guanaco} demonstrates unexpectedly robust Theory of Mind skills, as outlined in prior studies \cite{nematzadeh2018evaluating, sap2022neural}. For example, it gives detailed and accurate answers to prompts such as:

\begin{quote}
    \textbf{User}: James and Abby are in the bedroom. Abby put the pen in the desk drawer. Abby leaves the bedroom. James moves the pen into the bag. Where does James think Abby will look for the pen?

    \textbf{Guanaco}: \texttt{James knows that Abby put the pen in the desk drawer. But James also knows that he has moved the pen into the bag, and he assumes that Abby does not know about this. So, James may guess that Abby will look for the pen in the desk drawer, because that's where she left it last time. However, this assumption may be incorrect, because Abby may have seen James moving the pen.}
\end{quote}

Nonetheless, the model's reasoning in such contexts is inconsistent, and it often produces explanations that are incongruent with the actual scenario, as illustrated below:

\begin{quote}
    \textbf{User}: Evelyn entered the living room. Jackson entered the playroom. James entered the playroom. The beans are in the treasure chest. James exited the playroom. Jackson moved the beans to the pantry. Jackson exited the playroom. James entered the living room. Where will James look for the beans?

    \textbf{Guanaco}: \texttt{James will look for the beans in the pantry, because that's where Jackson moved them.}
\end{quote}

Here, {Guanaco} makes an unjustified assumption regarding knowledge transfer that was not described in the narrative. Such flaws reflect observations in contemporary research \cite{sap2022neural}, highlighting an area that merits deeper investigation.

\subsection{Considerations}
\label{ssec:considerations}

\paragraph{Evaluation}

Human annotators exhibited only moderate concordance (Fleiss $\kappa=0.42$), and this agreement further diminished when comparing outputs from two high-performing systems. These results expose shortcomings in existing benchmark datasets and evaluation methodologies for chatbots. In direct, manual assessments of responses generated by ChatGPT and {Guanaco} 65B using the Vicuna benchmark, we noticed that individual preferences among the paper's authors frequently diverged, indicating a strong subjective element. Advancing the field will require adopting or developing methodologies, potentially from Human-Computer Interaction or Psychology, to better accommodate subjective judgement.

Our investigation also reveals evident biases in automated assessment tools. Notably, GPT-4 tends to favor the output positioned first in its prompt, displaying pronounced order effects. Furthermore, the relatively low agreement between GPT-4 scores and those of human annotators (Fleiss $\kappa=0.25$) suggests a lack of alignment in evaluative criteria. Additionally, in Table~\ref{tbl:elo_large}, we find that GPT-4 awards substantially higher scores to its own responses versus those rated by humans, with Elo scores of 1348 compared to 1176—translating to an approximate 20\% increased chance of outperforming a peer.

Future research should assess potential biases present within automated evaluation frameworks, as well as explore strategies to address them.

\paragraph{Data \& Training} It is important to highlight that the OASST1 dataset, utilized for training the {Guanaco} models, encompasses multiple languages, and that the OA benchmark also incorporates prompts in various languages. Future research should investigate the extent to which multilingual training contributes to better performance on instructions issued in languages other than English, and whether this could account for the broader performance difference observed between the Vicuna-13B model (which was trained solely on English data) and the {Guanaco} 33B and 65B models evaluated on the OA benchmark.

Given the high performance demonstrated by the {Guanaco} models, we examined possible data leakage between OASST1 and the prompts found in the Vicuna benchmark. Our analyses, which included fuzzy string matching between the datasets and manual review of the closest pairs, revealed no shared prompts.

Additionally, it is worth mentioning that our models were trained exclusively via supervised learning using cross-entropy loss, without incorporating reinforcement learning from human feedback (RLHF). This underscores the need for more comprehensive examinations of the advantages and drawbacks associated with straightforward cross-entropy objectives in contrast to RLHF-based training. We anticipate that \textsc{QLoRA} will facilitate large-scale investigation of these issues, even in the absence of significant computational resources.

\section{Related Work}

\paragraph{Quantization of Large Language Models} Quantization approaches for large language models (LLMs) have predominantly emphasized optimizing inference performance. Leading strategies that preserve the performance of 16-bit LLMs target the mitigation of outlier activations (see, for instance, SmoothQuant~\citep{xiao2022smoothquant} and LLM.int8()~\citep{dettmers2022llm}), while alternative methods make use of advanced grouping schemes \citep{park2022nuqmm, yao2022zeroquant}. Studies on lossy quantization focus on evaluating trade-offs related to simple rounding mechanisms~\citep{dettmers2022case,zeng2022glm,pope2022efficiently} or propose techniques to optimize rounding decisions, thereby enhancing quantization accuracy~\citep{frantar2022gptq}. Apart from our contribution, only SwitchBack layers~\citep{wortsman2023stable} have investigated the application of backpropagation through quantized weights on models exceeding 1B parameters in scale.

\paragraph{Finetuning with Adapters} In our experiments, we leverage Low-rank Adapters~\citep{hu2021lora} (LoRA), yet the landscape of Parameter Efficient FineTuning (PEFT) strategies is broad, encompassing alternatives such as prompt tuning~\citep{qin2021learning,lester2021power,li2021prefix}, modifications to input embeddings~\citep{an2022input}, adaptation of intermediate representations (IA$^3$)~\citep{liu2022few}, incorporation of full additional layers~\citep{houlsby2019parameter}, adjustment of bias parameters~\citep{zaken2021bitfit}, applying learned masking over weights using Fisher information~\citep{sung2021training}, as well as hybrid methods~\citep{henderson2021compacter}. Our findings indicate that LoRA adapters can achieve parity with complete 16-bit finetuning. We suggest that future investigations systematically assess the comparative merits and tradeoffs among alternative PEFT methods.

\paragraph{Instruction Finetuning} To enhance a pretrained LLM’s capability to respond appropriately to prompts, instruction finetuning relies on datasets composed of input-output pairs sourced from various origins. This process further adapts a pretrained LLM to generate target outputs in response to given prompts. Diverse approaches and resources in this area include MetaICL~\citep{min2021metaicl}, MetaTuning~\citep{zhong2021adapting}, InstructGPT~\citep{ouyang2022training}, FLAN~\citep{wei2021finetuned,chung2022scaling}, PromptSource~\citep{bach2022promptsource}, Super-NaturalInstructions~\citep{wang2022super,sanh2021multitask}, Self-instruct~\citep{wang2022self}, UnnaturalInstructions~\citep{honovich2022unnatural}, OPT-IML~\citep{iyer2022opt}, UnifiedSKG~\citep{xie2022unifiedskg}, OIG/Chip2~\citep{laion2023}, Alpaca~\citep{alpaca}, Vicuna~\citep{vicuna2023}, Koala~\citep{koala_blogpost_2023}, and Self-instruct-GPT-4~\citep{peng2023instruction}.

\paragraph{Chatbots} Numerous models designed for instruction-following are deployed as chatbots, adopting conversational interaction patterns. Such systems frequently employ Reinforcement Learning from Human Feedback (RLHF)~\citep{christiano2017deep} or utilize training regimes that generate synthetic data from preexisting models, complemented by feedback from AI systems (RLAIF)~\citep{bai2022constitutional}. Pertinent works and datasets in this domain include Anthropic-HH~\citep{askell2021general,bai2022training}, Open Assistant~\citep{kopf2023openassistant}, LaMDA~\citep{thoppilan2022lamda}, and Sparrow~\citep{glaese2022improving}. In this study, we abstain from applying reinforcement learning; however, our top-performing model, Guanaco, is refined using the Open Assistant multi-turn chat dataset, which was curated for RLHF contexts~\citep{kopf2023openassistant}. For chatbot assessment, recent methodologies have replaced labor-intensive human annotation with automated evaluations using GPT-4~\citep{vicuna2023,peng2023instruction}. Our contribution includes improvements to these approaches, specifically prioritizing enhanced reliability in evaluation protocols.

\section{Limitations and Discussion}

We have demonstrated that our proposed approach, \textsc{QLoRA}, achieves performance equivalent to full 16-bit finetuning by utilizing a 4-bit foundation model together with Low-rank Adapters (LoRA). Nonetheless, we have yet to verify whether \textsc{QLoRA} can match 16-bit finetuning performance at the 33B and 65B parameter scales, due to the prohibitive computational demands; investigating these scales is deferred to subsequent research.

A further constraint relates to our assessment of instruction finetuned models. While evaluations are conducted using MMLU, the Vicuna benchmark, and the OA benchmark, our analysis does not include benchmarks such as BigBench, RAFT, or HELM, and there is no guarantee that our findings extrapolate to these datasets. Conversely, our work encompasses an extensive examination on MMLU and introduces novel chatbot evaluation techniques.

Based on the data discussed, the effectiveness of these benchmarks seems strongly related to the degree of overlap between the finetuning datasets and the evaluation benchmarks themselves. For instance, FLAN v2 shares substantial content with MMLU but not with conversational chatbot tests, while the Chip2 dataset aligns more with the latter; correspondingly, both models' performances reflect these similarities on the MMLU and Vicuna benchmarks, respectively. This underscores the importance of not only improving benchmark construction and evaluation protocols, but also of critically considering what abilities are actually being tested. Should our objective be to design models adept at recalling high school or collegiate factual knowledge, or to prioritize proficiency in conversational dialogue typical of chatbots? Or perhaps pursue other skills entirely? Given that established benchmarks are often favored due to the ease of evaluation over creating novel benchmarks, the selection of evaluation metrics may inadvertently guide research directions. It is therefore imperative for the research community to ensure that the chosen benchmarks truly capture the desired capabilities.

Although this work undertakes a comprehensive evaluation of overall chatbot capabilities, it remains limited in scope with respect to responsible AI assessments for {Guanaco}. Specifically, we compare the propensity for social bias in {Guanaco}-65B with that of other models in Table~\ref{tbl:crows}, finding that {Guanaco}-65B exhibits considerably lower average bias scores than unmodified pretrained models. These findings suggest that finetuning with the OASST1 dataset mitigates bias inherent in the original LLaMA model. Nevertheless, whether {Guanaco} shows reduced bias across a wider variety of contexts remains uncertain, and we defer more extensive bias analysis for {Guanaco} and analogous chatbots to forthcoming research.

A further constraint of this study is that it does not examine variations in bit precision—such as utilizing 3-bit models—or alternative adapter strategies. In addition to LoRA, numerous Parameter Efficient FineTuning (PEFT) methods have demonstrated effectiveness, but their scalability to larger models has yet to be systematically studied. LoRA was selected for its proven reliability, yet it remains possible that other adapters could offer superior results. Notably, post-quantization finetuning appears capable of recuperating much of the information lost during quantization, opening the possibility for more aggressive quantization approaches. As an example, it may be feasible that applying 3-bit GPTQ quantization to a basemodel, coupled with LoRA, could achieve performance on par with full 16-bit finetuning.

\section{Broader Impacts}

The \textsc{QLoRA} finetuning approach represents a significant advance by allowing 33B parameter models to be finetuned on a standard consumer GPU and 65B parameter models on a professional GPU, all while maintaining performance parity with full finetuning baselines. Our results indicate that the best-performing 33B model, trained on the Open Assistant dataset, matches ChatGPT on the Vicuna benchmark. Since instruction-based finetuning is a key mechanism for converting general pretrained LLMs into systems akin to ChatGPT, we anticipate that our technique will help democratize finetuning—especially benefiting researchers and practitioners with limited computational resources, and thus expanding access to state-of-the-art NLP tools. Consequently, \textsc{QLoRA} holds the potential to bridge resource gaps, reducing disparities between well-resourced organizations and smaller groups using consumer hardware.

An additional avenue for significant impact lies in the adaptation of the method for use on mobile devices. We contend that \textsc{QLoRA} could mark a pivotal step forward by making it feasible to finetune LLMs directly on smartphones and other environments with limited computational resources. Whereas previous efforts demonstrated that 7B parameter models can be executed on mobile hardware, \textsc{QLoRA} represents the first approach capable of supporting the finetuning process on such platforms. Our estimates suggest that, using an iPhone 12 Plus, it is possible to finetune approximately 3 million tokens overnight while the device is charging. Although finetuned 7B models do not attain the performance levels of systems like ChatGPT, their output quality appears sufficient to unlock new use cases that were previously restricted by either privacy constraints or subpar model performance. Through \textsc{QLoRA}, it becomes increasingly feasible to deploy LLMs in privacy-sensitive scenarios, empowering individuals to maintain ownership and governance over their data and models, while at the same time lowering the barriers to LLM deployment.

Nevertheless, it is important to acknowledge that finetuning also has dual-use potential, and could be leveraged for harmful purposes. The proliferation of LLMs comes with well-documented risks \citep{bommasani2021opportunities,bender2021dangers}. However, we maintain that democratizing access to this rapidly diffusing technology will enable more rigorous and independent examination, as opposed to consolidating the control of powerful LLMs within large corporations that restrict model and source code access and limit opportunities for external scrutiny.

In conclusion, we expect that \textsc{QLoRA} will facilitate broader and easier access to finetuning high-quality LLMs, thereby contributing to a positive societal impact.